\chapter{Conclusion}

\section{Summary}

This work set out to build a Software Transactional Memory system in Clojure
that closes a specific gap in the existing ecosystem: there was no safe way
to coordinate side-effects such as logging or notifications at transaction
boundaries. What came out of it is an STM implementation based on a variation
of TL2, extended with MVCC and a transactional event system, together with a
formal proof that every history the system produces satisfies
opacity~\cite{guerraoui2007}.

The evaluation results were not uniform, but they were informative.
Read-heavy workloads scale well, reaching 1.97M\,TPS at 12 threads, since
read-only transactions skip the commit lock entirely and run in genuine
parallel. High-contention workloads scale too, for a different reason: the
ownership protocol catches conflicts eagerly at write time rather than
waiting until commit. Write-heavy workloads are the exception: the global
commit lock serialises every write commit, so past the saturation point,
adding threads actively hurts throughput rather than helping it. As for the
transactional event system itself, it adds 0.22\,$\mu$s per commit, roughly
14\% over baseline, which is a small price given what it buys.

\section{Limitations and Future Work}
\label{sec:future}

\paragraph{Global commit lock.}
The most pressing limitation is the global commit lock. Section~\ref{sec:bottlenecks}
shows it accounts for nearly all thread blocking in write-heavy scenarios,
with total wait time growing faster than thread count once the lock
saturates. One fix stands out: replace the single lock with a fixed array of
per-stripe locks, hashing each reference to a stripe, following TL2's PS
policy~\cite{dice2006transactional}. Two write transactions landing on
different stripes could then commit in parallel, recovering the throughput
currently lost to serialisation without touching how an individual commit
works internally.

\paragraph{Dynamic version history sizing.}
The history-sweep results in Section~\ref{sec:history-sweep} show a fixed
global bound working well enough for the bank transfer workload, since trim
aborts never occur there, but nothing guarantees that holds more generally.
A workload combining short, frequent writes on a small number of references
with long-running read transactions could exhaust the history window before
a reader gets a chance to use it, producing trim aborts that a larger history
would have avoided. Simply raising \texttt{max-history} runs into the GC
overhead already visible at 12 threads for sizes 128 and 256. Clojure's own
STM avoids this trade-off by letting each reference track its own abort rate
and grow its history independently once trim aborts start occurring.
Something similar here would remove the need to hand-tune
\texttt{max-history}, and would likely cut abort rates in the workloads
where a fixed bound currently struggles.

\paragraph{Lifecycle hook granularity.}
The current hook model offers two attachment points, \texttt{on-commit} and
\texttt{on-rollback}, which covers most use cases but not all of them. A
\texttt{before-commit} hook, firing after read-set validation but before
writes are applied, would open the door to write-ahead logging: the pending
write set could be appended to a persistent log while the transaction still
holds the commit lock, giving durability guarantees the current model has no
way to express. Further phases such as \texttt{after-commit},
\texttt{before-retry}, and \texttt{after-retry} would extend this further,
covering use cases that need to act at specific points in the lifecycle
rather than only at its terminal boundaries.